%% ============================================================
%% cap02.tex — Capítulo 2
%% Requisitos del sistema de información
%% ============================================================

\chapter{Requisitos del sistema de información}
\label{cap:requisitos}

\noindent
Toda vez que el diseño orientado a objetos parte necesariamente de una
comprensión acuciosa del dominio del problema, el presente capítulo
se ocupa del análisis y la especificación de los requisitos del sistema
de información para el Estatuto General y el Estatuto Académico de la
\ud. El proceso de levantamiento de requisitos ---estructurado en los
Sprints~2 y 3--- no es un trámite burocrático previo al ``verdadero''
trabajo de programación: es, en sí mismo, un acto de modelado conceptual
que determina la calidad de las abstracciones que se construirán en los
capítulos subsiguientes.\footnote{%
  \textcite{Larman2004} sostiene que el análisis del dominio es el paso
  más crítico ---y con frecuencia el más descuidado--- en el desarrollo
  orientado a objetos. Un modelo de dominio deficiente propaga sus errores
  a lo largo de todas las fases del proyecto; los errores detectados en
  la fase de requisitos cuestan, en promedio, diez veces menos que los
  detectados en la fase de implementación \autocite{Pressman2014}.}
De suerte que la decisión de dedicar un capítulo completo a los requisitos
antes de escribir una sola línea de código expresa una convicción
pedagógica, no un protocolo rutinario.

%% ─────────────────────────────────────────────
\section{Análisis del dominio académico-administrativo}
\label{sec:analisis-dominio}
%% ─────────────────────────────────────────────

Se aplican técnicas de análisis del dominio ---glosarios, análisis de
documentos normativos, identificación de actores y procesos--- al corpus
compuesto por el \acuerdo, la propuesta de Estatuto Académico (noviembre
de 2025) y el documento de -PFA- de la \ud\ (2023). El objetivo es producir
un modelo conceptual preliminar que sirva de insumo para la especificación
formal de requisitos y, posteriormente, para el diagrama de clases del
dominio (Capítulo~\ref{cap:uml}).

\subsection{Identificación de actores y roles}
\label{subsec:actores-roles}

La Tabla~\ref{tab:actores-sistema} enumera los actores del sistema, con
referencia a los artículos del estatuto que los definen y a la necesidad
principal que cada uno tiene respecto del sistema de información.

\begin{longtable}{L{2.5cm} L{2.5cm} L{4cm} C{2cm}}
  \toprule
  \textbf{Actor} & \textbf{Rol en el sistema} &
  \textbf{Necesidad principal} & \textbf{Art.\ del estatuto} \\
  \midrule
  \endhead
  \bottomrule
  \caption{Actores del sistema de información y sus necesidades principales.
    Fuente: elaboración del autor con base en el \acuerdo.}
  \label{tab:actores-sistema}
  \endlastfoot
  Rector
  & Administrador institucional de máximo nivel
  & Consultar estado general de la institución; emitir resoluciones y acuerdos
  & Arts.\ 6, 18 \\
  \midrule
  Decano
  & Administrador de Facultad
  & Gestionar Facultad y sus dependencias; convocar Consejo de Facultad
  & Arts.\ 18, 49 \\
  \midrule
  Director de Escuela
  & Administrador de Escuela
  & Gestionar docentes adscritos; asignar carga académica; convocar Consejo de Escuela
  & Arts.\ 15, 18 \\
  \midrule
  Director de Programa
  & Coordinador académico de un programa
  & Gestionar plan de estudios; matricular estudiantes; gestionar evaluación y -PFA-
  & Arts.\ 18, 49 \\
  \midrule
  Coordinador de -CABA-
  & Coordinador de área de formación
  & Articular docencia e investigación en el área; gestionar actividades del -CABA-
  & Art.\ 19 \\
  \midrule
  Docente
  & Orientador de espacios académicos
  & Consultar asignación de espacios académicos; registrar evaluaciones y evidencias
  & Arts.\ 8, 15 \\
  \midrule
  Estudiante
  & Beneficiario del proceso formativo
  & Consultar matrícula; consultar calificaciones; acceder a información del programa
  & Arts.\ 8, 18 \\
  \midrule
  Personal administrativo
  & Operador del sistema en procesos administrativos
  & Gestionar matrículas, certificados, registros académicos
  & Art.\ 8 \\
  \midrule
  Egresado
  & Ex-miembro de la comunidad con derechos residuales
  & Consultar historial académico; solicitar certificaciones
  & Art.\ 8 \\
  \midrule
  Auditor interno
  & Actor externo de control
  & Consultar trazabilidad de decisiones de órganos colegiados; generar reportes de cumplimiento
  & (Implícito en el marco normativo) \\
\end{longtable}

\subsection{Procesos académicos derivados del Estatuto}
\label{subsec:procesos-academicos}

La Tabla~\ref{tab:procesos-academicos} mapea los procesos académico-administrativos
regulados por el estatuto que generan requisitos funcionales para el sistema,
con indicación de los actores involucrados y los datos que cada proceso
consume y produce.

\begin{longtable}{L{2.8cm} L{2.8cm} L{2.8cm} L{3cm}}
  \toprule
  \textbf{Proceso} & \textbf{Actores involucrados} &
  \textbf{Datos de entrada} & \textbf{Datos de salida} \\
  \midrule
  \endhead
  \bottomrule
  \caption{Procesos académicos del dominio con actores, entradas y salidas.}
  \label{tab:procesos-academicos}
  \endlastfoot
  Matrícula de estudiante
  & Estudiante, Personal administrativo, Director de Programa
  & Datos del aspirante; oferta académica vigente; requisitos del programa
  & Registro de matrícula; código de estudiante; estado ``activo'' \\
  \midrule
  Adscripción de docente a Escuela
  & Director de Escuela, Decano, Personal administrativo
  & Datos del docente; tipo de vinculación; Escuela de destino
  & Registro de adscripción; actualización del listado de la Escuela \\
  \midrule
  Asignación de espacios académicos
  & Director de Escuela, Docente
  & Plan de estudios vigente; disponibilidad docente; horarios
  & Asignación de docente a espacio académico; carga horaria \\
  \midrule
  Gestión del plan de estudios
  & Director de Programa, Consejo de Facultad
  & PEP vigente; espacios académicos; créditos; prerrequisitos
  & Plan de estudios actualizado; resolución de aprobación \\
  \midrule
  Registro de evaluación y -PFA-
  & Docente, Director de Programa
  & Rúbrica de evaluación; evidencias de aprendizaje; -PFA- del programa
  & Calificaciones registradas; informe de avance del estudiante frente a RA \\
  \midrule
  Sesión de órgano colegiado
  & Rector / Decano / Director (según órgano), representantes de la comunidad
  & Convocatoria; orden del día; quórum
  & Acta de sesión; decisiones (acuerdos, resoluciones) \\
  \midrule
  Creación de unidad académica
  & Rector, Consejo Superior (para Facultades); Decano, Consejo de Facultad
    (para Escuelas)
  & Estudio de factibilidad; resolución de creación
  & Nueva unidad registrada en el sistema; actualización de la jerarquía
    organizacional \\
\end{longtable}

%% ─────────────────────────────────────────────
\section{Requisitos funcionales}
\label{sec:req-funcionales}
%% ─────────────────────────────────────────────

Los requisitos funcionales del sistema se especifican en cinco grupos
correspondientes a los cinco subsistemas identificados en la
Sección~\ref{sec:caso-estudio}. La notación es RF-NNN, donde NNN
es un número secuencial de tres dígitos. Para cada requisito se indica
prioridad según el esquema MoSCoW: M (\textit{Must have}), S
(\textit{Should have}), C (\textit{Could have}).

\subsection{Gestión de estructura organizacional}
\label{subsec:req-organizacional}

\begin{longtable}{C{1.4cm} L{2.5cm} L{5cm} C{1.3cm} L{1.8cm}}
  \toprule
  \textbf{ID} & \textbf{Nombre} & \textbf{Descripción} &
  \textbf{Prior.} & \textbf{Subsistema} \\
  \midrule
  \endhead
  \bottomrule
  \caption{Requisitos funcionales: subsistema de gestión organizacional
    (RF-001 a RF-004). Fuente: \acuerdo, Arts.\ 14--19.}
  \label{tab:rf-organizacional}
  \endlastfoot
  RF-001
  & Registrar unidad académica
  & El sistema debe permitir registrar una nueva Facultad, Escuela, Centro,
    Instituto o CABA con sus atributos normativos (nombre, código, fecha de
    creación, director), validando que no exista duplicación de código
    y que la jerarquía organizacional sea coherente con el \acuerdo
  & M
  & SS-1 \\
  \midrule
  RF-002
  & Consultar estructura organizacional
  & El sistema debe permitir consultar la estructura completa de la \ud\
    de forma jerárquica (árbol Facultad $\rightarrow$ Escuela $\rightarrow$
    CABA), con filtros por nivel, nombre y estado (activa/inactiva)
  & M
  & SS-1 \\
  \midrule
  RF-003
  & Actualizar datos de unidad académica
  & El sistema debe permitir actualizar los atributos de una unidad
    académica existente, registrando la fecha y el actor responsable del
    cambio; las modificaciones deben dejar trazabilidad de auditoría
  & S
  & SS-1 \\
  \midrule
  RF-004
  & Gestionar directivos de unidades
  & El sistema debe permitir registrar y actualizar el directivo
    (decano, director de escuela, director de centro/instituto) de cada
    unidad académica, con validación de que el directivo sea un docente
    de planta adscrito a la institución
  & S
  & SS-1 \\
\end{longtable}

\subsection{Gestión de comunidad universitaria}
\label{subsec:req-comunidad}

\begin{longtable}{C{1.4cm} L{2.5cm} L{5cm} C{1.3cm} L{1.8cm}}
  \toprule
  \textbf{ID} & \textbf{Nombre} & \textbf{Descripción} &
  \textbf{Prior.} & \textbf{Subsistema} \\
  \midrule
  \endhead
  \bottomrule
  \caption{Requisitos funcionales: subsistema de gestión de comunidad
    universitaria (RF-005 a RF-008). Fuente: \acuerdo, Art.\ 8.}
  \label{tab:rf-comunidad}
  \endlastfoot
  RF-005
  & Registrar docente con tipo de vinculación
  & El sistema debe permitir registrar un docente con su tipo de vinculación
    (planta, ocasional, hora cátedra, visitante o experto), adscribiendo\-lo
    a exactamente una Escuela; los atributos requeridos difieren según
    el tipo de vinculación
  & M
  & SS-2 \\
  \midrule
  RF-006
  & Registrar y matricular estudiante
  & El sistema debe permitir registrar un estudiante y asociarlo a un
    programa académico mediante una matrícula; la matrícula registra el
    período académico, el estado y el número de créditos vigentes
  & M
  & SS-2 \\
  \midrule
  RF-007
  & Consultar comunidad universitaria
  & El sistema debe permitir consultar los miembros de la comunidad
    universitaria con filtros por tipo (docente, estudiante, personal
    administrativo, egresado), unidad académica y estado de vinculación
  & M
  & SS-2 \\
  \midrule
  RF-008
  & Gestionar egresados y certificaciones
  & El sistema debe permitir registrar el cambio de estado de un estudiante
    a egresado y generar certificaciones de grado, historia académica y
    créditos cursados, con firma digital del Director de Registro
  & S
  & SS-2 \\
\end{longtable}

\subsection{Gestión de programas académicos y planes de estudio}
\label{subsec:req-programas}

\begin{longtable}{C{1.4cm} L{2.5cm} L{5cm} C{1.3cm} L{1.8cm}}
  \toprule
  \textbf{ID} & \textbf{Nombre} & \textbf{Descripción} &
  \textbf{Prior.} & \textbf{Subsistema} \\
  \midrule
  \endhead
  \bottomrule
  \caption{Requisitos funcionales: subsistema de gestión académica
    (RF-009 a RF-012). Fuente: Estatuto Académico (propuesta), Arts.\ 13--18;
    \textcite{Decreto1330_2019}.}
  \label{tab:rf-academico}
  \endlastfoot
  RF-009
  & Registrar programa académico
  & El sistema debe permitir registrar un programa académico con su
    código -SNIES-, nivel (pregrado/posgrado), denominación, área de
    conocimiento, Facultad oferente y estado de registro calificado,
    conforme al Decreto~1330 de 2019
  & M
  & SS-3 \\
  \midrule
  RF-010
  & Gestionar plan de estudios
  & El sistema debe permitir crear y versionar el plan de estudios de
    un programa, con los espacios académicos (nombre, créditos, nivel
    de formación, prerrequisitos, corequisitos) y la distribución por
    semestres o períodos
  & M
  & SS-3 \\
  \midrule
  RF-011
  & Asignar docente a espacio académico
  & El sistema debe permitir asignar un docente a un espacio académico
    en un período específico, validando que el docente esté adscrito a
    la Escuela pertinente y que la carga horaria no supere los límites
    definidos por el tipo de vinculación
  & M
  & SS-3 \\
  \midrule
  RF-012
  & Gestionar PEP y documentos de política académica
  & El sistema debe permitir registrar y versionar el Proyecto Educativo
    de Programa (PEP), con referencia al PEF de la Facultad y al PUI
    de la institución; cada versión debe registrar fecha de aprobación
    y órgano que la aprobó
  & S
  & SS-3 \\
\end{longtable}

\subsection{Gestión de evaluación y resultados de aprendizaje}
\label{subsec:req-evaluacion}

\begin{longtable}{C{1.4cm} L{2.5cm} L{5cm} C{1.3cm} L{1.8cm}}
  \toprule
  \textbf{ID} & \textbf{Nombre} & \textbf{Descripción} &
  \textbf{Prior.} & \textbf{Subsistema} \\
  \midrule
  \endhead
  \bottomrule
  \caption{Requisitos funcionales: subsistema de gestión de evaluación y -PFA-
    (RF-013 a RF-016). Fuente: Documento -PFA- \ud\ (2023);
    \textcite{MEN2020_RA}.}
  \label{tab:rf-evaluacion}
  \endlastfoot
  RF-013
  & Definir Resultados de Aprendizaje por programa
  & El sistema debe permitir formular Resultados de Aprendizaje (RA) para
    un programa académico, con verbo de acción de la taxonomía de Bloom,
    condición, criterio de desempeño y vinculación al perfil de egreso;
    los RA deben ser observables y medibles \autocite{MEN2020_RA}
  & M
  & SS-4 \\
  \midrule
  RF-014
  & Registrar -PFA- por espacio académico
  & El sistema debe permitir registrar los Propósitos de Formación y de
    Aprendizaje (-PFA-) de un espacio académico, diferenciados por ámbito
    (ontológico, epistemológico, contextual), con vinculación a los RA
    del programa al que pertenece el espacio
  & M
  & SS-4 \\
  \midrule
  RF-015
  & Gestionar rúbricas de evaluación
  & El sistema debe permitir crear rúbricas de evaluación con criterios,
    niveles de desempeño (insuficiente, básico, alto, superior) y
    ponderaciones; cada rúbrica se asocia a uno o más -PFA- del espacio
    académico
  & M
  & SS-4 \\
  \midrule
  RF-016
  & Registrar evidencias y calcular avance
  & El sistema debe permitir al docente registrar evidencias de aprendizaje
    para un estudiante en un espacio académico y calcular automáticamente
    el nivel de avance del estudiante frente a los RA del programa;
    el cálculo debe ser trazable hasta los criterios de la rúbrica
  & S
  & SS-4 \\
\end{longtable}

\subsection{Gestión de órganos colegiados y decisiones}
\label{subsec:req-organos}

\begin{longtable}{C{1.4cm} L{2.5cm} L{5cm} C{1.3cm} L{1.8cm}}
  \toprule
  \textbf{ID} & \textbf{Nombre} & \textbf{Descripción} &
  \textbf{Prior.} & \textbf{Subsistema} \\
  \midrule
  \endhead
  \bottomrule
  \caption{Requisitos funcionales: subsistema de gestión de órganos colegiados
    (RF-017 a RF-020). Fuente: \acuerdo, Art.\ 18.}
  \label{tab:rf-organos}
  \endlastfoot
  RF-017
  & Registrar sesión de órgano colegiado
  & El sistema debe permitir registrar una sesión de cualquier órgano
    colegiado (Consejo Superior, Académico, de Facultad, Escuela, Centro
    o Instituto) con convocatoria, quórum, orden del día y lista de asistentes
  & M
  & SS-5 \\
  \midrule
  RF-018
  & Registrar decisiones (acuerdos y resoluciones)
  & El sistema debe permitir registrar las decisiones adoptadas en una
    sesión ---acuerdos y resoluciones, con número, fecha, extracto y
    clasificación temática--- y vincularlas a las unidades académicas
    o actores afectados
  & M
  & SS-5 \\
  \midrule
  RF-019
  & Consultar historial de decisiones
  & El sistema debe permitir consultar el historial de decisiones de
    cualquier órgano colegiado, con filtros por fecha, tipo de decisión,
    unidad afectada y estado (vigente/derogado)
  & S
  & SS-5 \\
  \midrule
  RF-020
  & Generar reportes de cumplimiento normativo
  & El sistema debe permitir generar reportes que evidencien el cumplimiento
    de las disposiciones del \acuerdo\ y del Estatuto Académico, cruzando
    las decisiones registradas con las obligaciones normativas derivadas
    de cada artículo; útil para procesos de autoevaluación y renovación
    de registro calificado
  & C
  & SS-5 \\
\end{longtable}

%% ─────────────────────────────────────────────
\section{Requisitos no funcionales}
\label{sec:req-no-funcionales}
%% ─────────────────────────────────────────────

Los requisitos no funcionales (-RNF-) establecen las condiciones de calidad
bajo las cuales el sistema debe operar. En la medida en que el sistema
implementará procesos normativos de una institución pública de educación
superior, los requisitos de trazabilidad, seguridad y mantenibilidad
adquieren carácter determinante. La Tabla~\ref{tab:rnf} los especifica.\footnote{%
  La clasificación ISO/IEC 25010 (Quality Model for Software Product) se
  emplea como referencia para las categorías de los -RNF-. Sin embargo,
  la clasificación se adapta al dominio académico, priorizando las
  categorías con mayor incidencia en las decisiones de diseño orientado
  a objetos.}

\begin{longtable}{C{1.6cm} L{1.8cm} L{5.5cm} L{3cm}}
  \toprule
  \textbf{ID} & \textbf{Categoría} & \textbf{Descripción} &
  \textbf{Métrica de verificación} \\
  \midrule
  \endhead
  \bottomrule
  \caption{Requisitos no funcionales del sistema (RNF-001 a RNF-008).
    Categorías según ISO/IEC 25010.}
  \label{tab:rnf}
  \endlastfoot
  RNF-001
  & Rendimiento
  & Las consultas sobre la estructura organizacional y los datos del
    estudiante deben responder en tiempo inferior a 2 segundos para
    hasta 5.000 usuarios concurrentes
  & Prueba de carga con JMeter; p95 $\leq$ 2 s bajo 5.000 usuarios concurrentes \\
  \midrule
  RNF-002
  & Seguridad (autenticación)
  & El sistema debe implementar autenticación multifactor para los roles
    con acceso a datos sensibles (rector, decano, director de programa,
    personal administrativo); los docentes y estudiantes pueden autenticarse
    con usuario y contraseña
  & 100\% de usuarios con acceso a datos personales requieren MFA;
    auditado en prueba de penetración \\
  \midrule
  RNF-003
  & Seguridad (integridad de datos)
  & Toda modificación de datos debe registrar al actor responsable,
    la fecha y hora, y el valor anterior; los registros de auditoría
    no pueden ser modificados ni eliminados por ningún actor del sistema
  & Revisión de trazas de auditoría; prueba de intento de modificación
    no autorizada \\
  \midrule
  RNF-004
  & Usabilidad
  & El sistema debe ser usable por personal administrativo sin formación
    técnica especializada; el tiempo de aprendizaje para las tareas
    más frecuentes no debe superar 2 horas
  & Prueba de usabilidad con 5 usuarios no técnicos; tasa de error $\leq$ 10\%
    en las 5 tareas más frecuentes \\
  \midrule
  RNF-005
  & Portabilidad
  & El sistema debe ejecutarse sin modificación de código en entornos
    \textit{Linux} (Debian/Ubuntu) y \textit{Windows} 10/11; la versión
    \java\ debe ser compatible con JDK 17 LTS o superior
  & Ejecución de la suite de pruebas en ambos entornos sin fallos;
    cero dependencias de sistema operativo específico \\
  \midrule
  RNF-006
  & Mantenibilidad
  & El código fuente debe cumplir los principios SOLID y lograr una
    cobertura de pruebas unitarias superior al 80\%; la deuda técnica
    medida con SonarQube no debe superar 5\% en el nivel ``crítico''
  & Reporte SonarQube; métrica de cobertura con JaCoCo (Java) o Cargo tarpaulin (\rust) \\
  \midrule
  RNF-007
  & Disponibilidad
  & El sistema debe estar disponible el 99\% del tiempo en horario
    hábil de la institución (lunes a viernes, 7:00--22:00); las
    ventanas de mantenimiento se programarán en horario no hábil
  & Monitoreo con herramienta de disponibilidad; registro de incidentes \\
  \midrule
  RNF-008
  & Trazabilidad normativa
  & Cada requisito funcional implementado debe estar vinculado a al menos
    un artículo del \acuerdo\ o del Estatuto Académico; la trazabilidad
    debe ser verificable mediante una matriz RF $\leftrightarrow$ artículo
  & Matriz de trazabilidad RF $\leftrightarrow$ estatuto; revisión por pares
    antes de cada entrega \\
\end{longtable}

%% ─────────────────────────────────────────────
\section{Historias de usuario iniciales}
\label{sec:historias-usuario}
%% ─────────────────────────────────────────────

Las historias de usuario del \textit{product backlog} inicial se formulan
con la plantilla estándar: «Como [actor], quiero [funcionalidad], para que
[valor de negocio]». Cada historia incluye criterios de aceptación
verificables.\footnote{%
  La elección del formato de historias de usuario, en lugar de casos de uso
  formales de -UML-, obedece a una decisión pedagógica: los estudiantes de
  primer semestre de la asignatura encuentran más accesible el lenguaje
  natural de las historias. Los casos de uso formales se desarrollan en el
  Capítulo~\ref{cap:uml} a partir de estas historias.}

\bigskip

\textbf{HU-001.} Como \textbf{Director de Escuela}, quiero registrar la
adscripción de un nuevo docente a mi escuela, para que el sistema refleje
la composición actualizada del cuerpo docente y la asignación de carga
académica sea coherente con la planta disponible.

\textit{Criterios de aceptación:}
\begin{itemize}
  \item El sistema valida que el tipo de vinculación sea uno de los cinco
    definidos en el Art.\ 8 del \acuerdo\ (planta, ocasional, hora cátedra,
    visitante, experto).
  \item El sistema verifica que el docente no esté ya adscrito a otra
    Escuela de la misma Facultad.
  \item La adscripción queda registrada con fecha, actor responsable y
    referencia al acto administrativo que la respalda.
  \item El Director de Escuela recibe confirmación inmediata y puede
    consultar la nueva nómina actualizada.
\end{itemize}

\bigskip

\textbf{HU-002.} Como \textbf{Estudiante}, quiero consultar mi plan
de estudios vigente con el estado de avance en cada espacio académico,
para que pueda planificar los períodos académicos restantes para completar
el programa.

\textit{Criterios de aceptación:}
\begin{itemize}
  \item El sistema muestra el plan de estudios completo con semestres,
    espacios académicos, créditos, prerrequisitos y estado de cada espacio
    (pendiente, en curso, aprobado, reprobado).
  \item El porcentaje de avance se calcula sobre el total de créditos del
    programa.
  \item La información está disponible en menos de 2 segundos (RNF-001).
  \item El estudiante no puede modificar su plan de estudios mediante
    esta funcionalidad; solo puede consultarlo.
\end{itemize}

\bigskip

\textbf{HU-003.} Como \textbf{Docente}, quiero registrar las evidencias
de aprendizaje de un estudiante en mi espacio académico con la rúbrica
asociada, para que el seguimiento del avance del estudiante frente a los
Resultados de Aprendizaje del programa sea trazable.

\textit{Criterios de aceptación:}
\begin{itemize}
  \item El sistema presenta la rúbrica del espacio académico con sus criterios
    y niveles de desempeño.
  \item El docente puede seleccionar el nivel de desempeño para cada criterio
    y adjuntar un archivo de evidencia (máx.\ 20 MB).
  \item El sistema calcula automáticamente la calificación según la
    ponderación de la rúbrica.
  \item El registro queda sellado con fecha y hora; el docente puede modificarlo
    hasta 48 horas después del registro inicial.
\end{itemize}

\bigskip

\textbf{HU-004.} Como \textbf{Decano}, quiero consultar la estructura
organizacional completa de mi Facultad ---Escuelas, Centros, Institutos
y CABAS--- con el número de docentes adscritos a cada Escuela y el número
de programas activos, para que pueda tomar decisiones informadas sobre
la distribución de recursos humanos y académicos.

\textit{Criterios de aceptación:}
\begin{itemize}
  \item El reporte muestra la jerarquía completa de la Facultad hasta el
    nivel de CABA.
  \item Para cada Escuela, muestra el número de docentes por tipo de vinculación.
  \item Para cada programa activo, muestra el número de estudiantes matriculados
    en el período vigente.
  \item El reporte puede exportarse en formato PDF.
\end{itemize}

\bigskip

\textbf{HU-005.} Como \textbf{Secretario de Consejo de Facultad}, quiero
registrar los acuerdos adoptados en una sesión del Consejo de Facultad,
para que las decisiones queden documentadas con trazabilidad normativa
y puedan consultarse en procesos de autoevaluación.

\textit{Criterios de aceptación:}
\begin{itemize}
  \item El registro incluye número de acuerdo, fecha, asunto, texto del
    articulado y referencia a las entidades o personas afectadas.
  \item Los registros de auditoría no pueden modificarse una vez guardados
    (RNF-003).
  \item El sistema vincula cada acuerdo a los artículos del \acuerdo\ que
    lo fundamentan.
  \item El listado de acuerdos puede consultarse con filtros por fecha,
    tipo y estado (vigente/derogado).
\end{itemize}

\bigskip

\textbf{HU-006.} Como \textbf{Director de Programa}, quiero registrar
los Propósitos de Formación y de Aprendizaje (-PFA-) de cada espacio
académico del programa, vinculados a los Resultados de Aprendizaje del
plan de estudios, para que la evaluación formativa sea coherente con el
proyecto educativo del programa.

\textit{Criterios de aceptación:}
\begin{itemize}
  \item El sistema permite especificar -PFA- en los tres ámbitos (ontológico,
    epistemológico, contextual) para cada espacio académico.
  \item Cada -PFA- se vincula a al menos un Resultado de Aprendizaje del programa.
  \item El sistema verifica que todos los RA del programa tengan al menos un
    -PFA- vinculado en algún espacio del plan de estudios.
  \item La configuración de -PFA- puede versionarse al inicio de cada período
    académico.
\end{itemize}

\bigskip

\textbf{HU-007.} Como \textbf{Personal administrativo}, quiero registrar
la matrícula de un nuevo estudiante en un programa académico para el período
vigente, para que el sistema académico refleje correctamente la población
estudiantil activa y los cupos disponibles.

\textit{Criterios de aceptación:}
\begin{itemize}
  \item El sistema verifica que el programa tenga cupos disponibles en el
    período seleccionado.
  \item La matrícula incluye: estudiante, programa, período, modalidad,
    créditos inscritos y valor de matrícula.
  \item El sistema genera automáticamente un código de matrícula único.
  \item El estudiante queda con estado ``activo'' y puede consultar su
    información académica.
\end{itemize}

\bigskip

\textbf{HU-008.} Como \textbf{Rector}, quiero consultar un tablero de
control con indicadores institucionales clave ---número de programas activos
con registro calificado vigente, número de docentes por tipo de vinculación,
número de estudiantes matriculados por Facultad, número de decisiones de
órganos colegiados en el período--- para que la toma de decisiones
estratégicas se sustente en datos actualizados.

\textit{Criterios de aceptación:}
\begin{itemize}
  \item El tablero se actualiza en tiempo real a partir de los datos del
    sistema.
  \item Cada indicador puede desagregarse por Facultad, período académico
    y tipo de vinculación.
  \item El tablero puede compartirse en modo de solo lectura con el
    Consejo Superior.
  \item Los datos del tablero coinciden con los reportes de los subsistemas
    individuales.
\end{itemize}

%% ─────────────────────────────────────────────
\section{Modelo de contexto del sistema}
\label{sec:modelo-contexto}
%% ─────────────────────────────────────────────

El diagrama de contexto establece los límites del sistema de información
frente a los sistemas externos con los que interactúa y los actores
que se relacionan directamente con él. Se representa en \plantuml\ como
diagrama de casos de uso de alto nivel:

\begin{lstlisting}[style=plantumlStyle,
  caption={Diagrama de contexto del sistema de información de la \ud\
    (PlantUML — nivel de contexto).},
  label={lst:plantuml-contexto}]
@startuml contexto_sistema
skinparam actorStyle awesome
skinparam rectangleBackgroundColor #F0F4FF
skinparam arrowColor #00467F

actor "Rector / Vicerrector" as rector
actor "Decano" as decano
actor "Director de Escuela" as director_escuela
actor "Director de Programa" as director_programa
actor "Docente" as docente
actor "Estudiante" as estudiante
actor "Personal administrativo" as admin

rectangle "Sistema de Información UD\n(Estatuto General y Académico)" as SI {
}

rectangle "Sistemas externos" {
  rectangle "Sistema de Nómina" as nomina
  rectangle "MEN / SNIES" as snies
  rectangle "Plataforma e-Learning" as elearning
  rectangle "Sistema de Biblioteca" as biblioteca
  rectangle "Sistema de Bienestar" as bienestar
}

rector --> SI : gestión institucional
decano --> SI : gestión de Facultad
director_escuela --> SI : gestión de Escuela y docentes
director_programa --> SI : gestión académica y PFA
docente --> SI : registro de evaluaciones
estudiante --> SI : consulta académica
admin --> SI : operación de registros

SI --> nomina : datos de vinculación docente
SI <-- nomina : confirmación de nómina
SI --> snies : reporte de programas activos
SI <-- elearning : actividades y evidencias
SI --> biblioteca : solicitudes de recursos
SI --> bienestar : datos de estudiantes en riesgo

@enduml
\end{lstlisting}

Los flujos de información principales entre el sistema y los actores
externos merecen una descripción sucinta. La integración con el sistema
de nómina es de carácter bidireccional: el sistema de información
reporta los datos de vinculación de los docentes y recibe confirmación
del estado de nómina. La integración con el -SNIES- del Ministerio de
Educación Nacional es unidireccional: el sistema genera los reportes
de programas activos, estudiantes matriculados y egresados que la
normativa colombiana obliga a reportar \autocite{Decreto1330_2019}.
La plataforma de \textit{e-learning} ---cuya integración podría
implementarse mediante una -API- -REST- en una fase posterior--- aportaría
las evidencias de aprendizaje registradas en entornos virtuales, de suerte
que el sistema de evaluación contaría con datos más completos para el
cálculo del avance del estudiante frente a los RA. Queda por determinar
si esta integración es prioritaria para una primera versión del sistema.

%% ─────────────────────────────────────────────
\section{Comparativa \java\ vs.\ \rust: impacto de los requisitos no funcionales}
\label{sec:java-rust-requisitos}
%% ─────────────────────────────────────────────

Los requisitos no funcionales especificados en la Sección~\ref{sec:req-no-funcionales}
condicionan de manera diferenciada las decisiones de implementación en
\java\ y en \rust. Tres dimensiones resultan particularmente reveladoras.\footnote{%
  Queda por determinar si la comparación entre \java\ y \rust\ desde la
  perspectiva de los -RNF- es generalmente válida o si depende
  excesivamente del dominio específico del sistema. En dominios con alta
  concurrencia y requisitos estrictos de latencia, la ventaja de \rust\
  sería más pronunciada; en dominios con predominio de operaciones
  de entrada/salida y acceso a base de datos, la diferencia podría ser
  menos significativa. La extrapolación exige cautela.}

\textbf{Seguridad de memoria (RNF-002, RNF-003).} El requisito de
integridad de datos y la imposibilidad de modificar registros de auditoría
resultan más fáciles de garantizar en \rust\, toda vez que el sistema
de \textit{ownership} impide la modificación de datos compartidos sin
paso explícito por \texttt{Mutex} o \texttt{RwLock}. En \java\, la
mutabilidad implícita de los objetos exige disciplina adicional: el
uso de colecciones inmutables (\texttt{Collections.unmodifiableList})
y de objetos de valor (\textit{value objects}) para las entidades
que no deben modificarse una vez creadas \autocite{Bloch2018}.

\textbf{Concurrencia (RNF-001, RNF-007).} El requisito de responder a
5.000 usuarios concurrentes en menos de 2 segundos plantea retos de
concurrencia que ambos lenguajes abordan de manera disímil. En \java\,
el \textit{framework} \texttt{java.util.concurrent} proporciona
mecanismos maduros y documentados; en \rust\, la garantía de ausencia
de condiciones de carrera (\textit{fearless concurrency}) reduce
el riesgo de errores difíciles de detectar \autocite{Klabnik2023}.

\textbf{Mantenibilidad (RNF-006).} El requisito de cobertura superior
al 80\% y cumplimiento de principios SOLID es independiente del lenguaje,
pero el ecosistema de herramientas de análisis estático es más maduro
en \java\ ---SonarQube, SpotBugs, PMD--- que en \rust\ ---Clippy, Miri---,
lo que podría facilitar el cumplimiento del -RNF- durante el desarrollo.

%% ─────────────────────────────────────────────
%% SPRINT 2
%% ─────────────────────────────────────────────
\section*{Sprint 2: Levantamiento de requisitos}
\label{sprint:2}
\addcontentsline{toc}{section}{Sprint 2: Levantamiento de requisitos}

\begin{tcolorbox}[sprintBox, title={Sprint 2 — Levantamiento de requisitos}]
  \textbf{Duración estimada:} 1 semana\\
  \textbf{Objetivo:} Especificar los requisitos funcionales y no funcionales
  del sistema a partir del análisis del \acuerdo\ y del Estatuto Académico,
  y formular las historias de usuario del \textit{product backlog} inicial.\\
  \textbf{Entregables:} (1) Documento RF (mínimo 20) con referencia normativa;
  (2) Documento RNF (mínimo 6); (3) \textit{Product backlog} con 8 historias
  priorizadas.
\end{tcolorbox}

\bigskip

\subsection*{Tabla del Sprint 2}

\begin{longtable}{L{3.5cm} L{4.5cm} L{4.5cm}}
  \sprintcabecera
  \endhead
  \bottomrule
  \caption{Entregables, características del programa y criterios de la
    rúbrica del Sprint 2.}
  \label{tab:sprint2}
  \endlastfoot
  Documento de requisitos funcionales (mínimo 20 RF, agrupados por
  subsistema, con referencia al artículo del estatuto, prioridad
  MoSCoW y subsistema)
  & Aplicar técnicas de levantamiento de requisitos a un dominio
    normativo real; distinguir actores, procesos y datos del dominio
  & El estudiante especifica al menos 20 RF coherentes con el estatuto,
    correctamente agrupados, con prioridad justificada y referencia
    normativa explícita a los artículos del \acuerdo\ o del Estatuto
    Académico \\
  \midrule
  Documento de requisitos no funcionales (mínimo 6 RNF, con categoría,
  descripción y métrica de verificación)
  & Identificar restricciones de calidad del sistema derivadas del
    dominio normativo e institucional
  & Los RNF son verificables (tienen métricas concretas), coherentes con
    el dominio institucional y relevantes para las decisiones de
    implementación en \java\ y \rust \\
  \midrule
  \textit{Product backlog} inicial (mínimo 8 historias de usuario
  priorizadas con MoSCoW, con criterios de aceptación para las
  3 de mayor prioridad)
  & Formular historias de usuario centradas en los actores del dominio
    institucional; articular requisitos funcionales con necesidades
    reales de los usuarios
  & Las historias identifican correctamente el actor, la funcionalidad
    y el valor; los criterios de aceptación son verificables y coherentes
    con el dominio institucional \\
  \midrule
  Matriz de trazabilidad RF $\leftrightarrow$ artículo del estatuto
  (tabla con al menos 20 filas)
  & Evidenciar la trazabilidad entre requisitos y marco normativo;
    aplicar el principio de trazabilidad normativa (RNF-008)
  & Cada RF está vinculado a al menos un artículo del \acuerdo\ o del
    Estatuto Académico; la matriz es verificable mediante consulta
    directa al texto del estatuto \\
\end{longtable}

\subsection*{Actividad de reflexión: software generativo y levantamiento de requisitos}
\label{act:ia-sprint2}

\begin{tcolorbox}[actividadIA,
  title={Actividad de reflexión (Sprint~2) —
         Software generativo y levantamiento de requisitos}]

\textbf{Instrucción general.}
El estudiante elabora un documento de requisitos funcionales a partir
de la lectura directa del \acuerdo\ y del Estatuto Académico, sin
asistencia computacional. A continuación, solicita a un asistente
generativo que derive los mismos requisitos a partir de los mismos
documentos (proporcionados como contexto). Se comparan ambas propuestas
aplicando el marco de análisis crítico de la Sección~\ref{sec:ia-reflexion-intro}.

\bigskip
\textbf{Preguntas orientadoras:}
\begin{enumerate}[label=\arabic*.]
  \item ¿El asistente omite requisitos que usted identificó? ¿A qué
        artículos del estatuto corresponden?
  \item ¿El asistente genera requisitos que no tienen respaldo en el
        estatuto (categorizables como ``ruido'' o ``inventados'')? ¿Cómo
        los detectó?
  \item ¿En qué medida el asistente distingue entre requisitos funcionales
        y no funcionales? ¿La distinción es coherente con la que se
        establece en la sección~\ref{sec:req-no-funcionales}?
  \item ¿El asistente genera RF relacionados con subsistemas que usted
        no había considerado? ¿Son pertinentes?
\end{enumerate}

\bigskip
\textbf{Producto esperado:} Informe de análisis comparativo (máximo
4 páginas) con tabla de requisitos propios/asistente, diferencias
detectadas y conclusiones argumentadas.

\bigskip
\textbf{Rúbrica breve:}
\begin{itemize}
  \item \textbf{Producto} (60\%): completitud ($\geq$ 20 RF propios),
        coherencia normativa, corrección de la especificación, métricas
        verificables en los RNF.
  \item \textbf{Proceso} (40\%): calidad del análisis crítico,
        profundidad de la comparación, identificación de errores del
        asistente con argumentación normativa específica.
\end{itemize}

\end{tcolorbox}

%% ─────────────────────────────────────────────
%% SPRINT 3
%% ─────────────────────────────────────────────
\section*{Sprint 3: Validación de requisitos con \textit{stakeholders}}
\label{sprint:3}
\addcontentsline{toc}{section}{Sprint 3: Validación de requisitos con
  \textit{stakeholders}}

\begin{tcolorbox}[sprintBox, title={Sprint 3 — Validación de requisitos con
  \textit{stakeholders}}]
  \textbf{Duración estimada:} 1 semana\\
  \textbf{Objetivo:} Validar los requisitos especificados en el Sprint~2
  mediante técnicas de revisión por pares y entrevistas simuladas con
  actores del dominio institucional; refinar el \textit{product backlog}.\\
  \textbf{Entregables:} (1) Acta de entrevista de validación; (2)
  \textit{Backlog} refinado con criterios de aceptación para las 5
  historias de mayor prioridad; (3) Informe de reflexión crítica.
\end{tcolorbox}

\bigskip

\subsection*{Tabla del Sprint 3}

\begin{longtable}{L{3.5cm} L{4.5cm} L{4.5cm}}
  \sprintcabecera
  \endhead
  \bottomrule
  \caption{Entregables, características del programa y criterios de la
    rúbrica del Sprint 3.}
  \label{tab:sprint3}
  \endlastfoot
  Acta de entrevista de validación con al menos un actor del dominio
  (docente, director de programa o personal administrativo de la \ud;
  o simulado mediante \textit{roleplay} documentado)
  & Aplicar técnicas de validación de requisitos con actores reales
    o simulados del dominio institucional; identificar ambigüedades,
    omisiones y contradicciones en la especificación
  & El acta registra con fidelidad los aportes del actor entrevistado,
    los contrasta con los requisitos especificados e identifica al
    menos 3 ajustes necesarios con justificación \\
  \midrule
  \textit{Product backlog} refinado con criterios de aceptación
  para las 5 historias de mayor prioridad (mínimo 3 criterios
  por historia, formulados en formato Given-When-Then)
  & Formular criterios de aceptación verificables y precisos para
    las historias de usuario prioritarias; aplicar el formato
    Given-When-Then
  & Los criterios de aceptación son observables, medibles, coherentes
    con el dominio institucional y expresados en formato
    Given-When-Then \\
  \midrule
  Informe de análisis crítico de los criterios de aceptación
  generados por asistente computacional frente a los criterios
  propios (ver actividad de reflexión)
  & Desarrollar pensamiento crítico frente a los criterios generados
    automáticamente; identificar limitaciones del asistente en el
    tratamiento de dominios normativos institucionales específicos
  & El informe identifica al menos 3 diferencias sustanciales entre
    los criterios propios y los del asistente, con argumentación
    basada en el dominio institucional y en los artículos del \acuerdo \\
  \midrule
  Diagrama de contexto refinado (actualización del diagrama inicial
  del Sprint~1 con base en los hallazgos de la validación)
  & Incorporar retroalimentación de la validación al modelo de contexto;
    refinar los flujos de información con actores reales
  & El diagrama actualizado refleja fielmente los hallazgos de la
    validación; los cambios respecto de la versión inicial están
    justificados en el acta \\
\end{longtable}

\subsection*{Actividad de reflexión: software generativo y validación de requisitos}
\label{act:ia-sprint3}

\begin{tcolorbox}[actividadIA,
  title={Actividad de reflexión (Sprint~3) —
         Software generativo y validación de requisitos}]

\textbf{Instrucción general.}
El estudiante utiliza un asistente computacional generativo para generar
criterios de aceptación para las 5 historias de usuario de mayor prioridad
del \textit{backlog} refinado. Se evalúa la pertinencia y la especificidad
institucional de los criterios generados en contraste con los criterios
propios derivados del proceso de validación con actores reales o simulados.

\bigskip
\textbf{Preguntas orientadoras:}
\begin{enumerate}[label=\arabic*.]
  \item ¿Los criterios de aceptación generados por el asistente son
        verificables en el contexto real de la \ud? Proporcione al menos
        dos ejemplos concretos de criterios que requieran ajuste.
  \item ¿El asistente produce criterios genéricos que podrían aplicarse
        a cualquier institución universitaria sin modificación? ¿Cómo
        se evidencia la falta de especificidad institucional?
  \item ¿Qué fragmentos del \acuerdo\ o del Estatuto Académico fueron
        necesarios proporcionar al asistente para obtener criterios más
        pertinentes? ¿La calidad mejoró con el contexto adicional?
  \item Si el asistente generó criterios de aceptación que usted no había
        considerado y que resultan válidos, ¿cómo integraría este hallazgo
        en su reflexión sobre el uso responsable del software generativo?
\end{enumerate}

\bigskip
\textbf{Rúbrica breve:}
\begin{itemize}
  \item \textbf{Producto} (60\%): pertinencia institucional de los
        criterios de aceptación propios, verificabilidad, uso correcto
        del formato Given-When-Then, completitud de la cobertura de las
        5 historias prioritarias.
  \item \textbf{Proceso} (40\%): calidad del análisis crítico,
        identificación argumentada de limitaciones del asistente,
        evidencia de iteración entre criterios propios y retroalimentación
        del asistente.
\end{itemize}

\end{tcolorbox}
